\section{Imagining Biophilic Interaction in Indoor Environments}

Having examined how an affective lens can help HBI understand existing indoor experience, the dissertation shifts towards design. This transition raises a further question: how can experiential knowledge inform future indoor environments without being converted into requirements for producing predetermined emotional states?

To address this question, the thesis engages architecture, spatial-design, and related experts. Their perspectives are important because the design of indoor environments involves knowledge of spatial organisation, sensory conditions, materiality, atmosphere, building use, and architectural practice. The expert studies bring these forms of knowledge into dialogue with the experiential concerns of HCI and HBI.

The first study consists of semi-structured interviews examining expert imaginaries of multisensory biophilic design.\sidenote{I use \emph{imaginaries} to refer to future-oriented accounts of what indoor environments could become. The term concerns design possibility rather than prediction.} Rather than asking experts to evaluate an established system, the interviews invite them to articulate possibilities, tensions, and directions for future indoor environments.

Biophilic design provides a productive context for this inquiry because it concerns how relations with nature can be incorporated into built environments. It is frequently represented through plants, daylight, natural materials, water, views, or organic forms. These strategies are important, but they do not encompass the full experiential potential of relations with nature.

Nature-related experience can also involve sound, movement, temporal variation, atmosphere, ambiguity, sensory correspondence, and changing environmental conditions. These qualities connect to the affective account developed through the occupant studies. They provide ways of reasoning about experience without assuming that design should generate a single desirable emotional state.

The interviews identify dimensions of multisensory biophilic design that remain under-explored within technologically mediated indoor environments. They direct attention beyond the visual presence of nature towards environmental processes, sensory relationships, social experience, spatial context, and the ways interactive technologies may support or complicate these relations.

The interviews establish the relevant dimensions, but they do not fully demonstrate how experts would express them through design. The dissertation therefore proceeds to expert design sessions in which participants articulate, visualise, and discuss possible multisensory biophilic environments. These sessions examine how future-oriented ideas are translated into spatial and interactive propositions.

Together, the interviews and design sessions produce a biophilic design framework for HBI, with broader implications for HCI. The framework organises the expert findings into design approaches for technologically mediated indoor environments. It does not prescribe a fixed set of natural elements or system features. Its purpose is to expand the design vocabulary of HBI by connecting affective, multisensory, spatial, and nature-related experience.
